% SERAPIS v3 — draft. Generated section by section by the readiness ladder (R5). AI-assisted; see Acknowledgements.
\documentclass[aps,prd,twocolumn,superscriptaddress,nofootinbib]{revtex4-2}
\usepackage{amsmath,amssymb,graphicx,hyperref}
\begin{document}
\title{SERAPIS: a modular plateau with a massive vacuum at the square torus, its Janus pair of predictions, and a three-observable consistency relation}
\author{Danny Eldridge}
\affiliation{Independent researcher\\ [AFFILIATION / CORRESPONDENCE EMAIL: Danny to confirm]}
\date{Draft of 15 November 2026}
\begin{abstract}
[ABSTRACT: written in R7.]
\end{abstract}
\maketitle

%==SEC1==
\section{Introduction}
Plateau models of inflation, and in particular $\alpha$-attractors, predict a scalar tilt close to $n_s \simeq 1-2/N_*$ and a small tensor-to-scalar ratio $r$. The latest combined analyses move the preferred tilt upward: CMB data (Planck, ACT DR6, SPT-3G with BICEP/Keck) give $n_s = 0.9682 \pm 0.0032$, and adding DESI DR2 baryon acoustic oscillations gives $n_s = 0.9728 \pm 0.0029$, with $r < 0.035$ at 95\% confidence~\cite{arXiv_2512_10613}. A plain plateau at $N_* \approx 50$ sits about $3$--$4\sigma$ below the second value.

Modular-invariant versions of these models, in which the inflaton is the shape modulus $\tau$ of a two-torus and the potential is invariant under $SL(2,\mathbb{Z})$, have been developed in detail~\cite{arXiv_2408_05203,arXiv_2412_19324,arXiv_2405_06497,arXiv_2407_12081,arXiv_2411_18603}. They tie inflation to a geometric question: at which shape does the torus come to rest after inflation?

In this paper we study a model, which we call SERAPIS, whose vacuum sits at the square torus $\tau = i$ (and a variant at the hexagonal torus $\rho = e^{2\pi i/3}$). We make four points.
(i) Symmetry alone forces $\tau = i$ to be a critical point and a natural zero of the potential, but a plateau built from $|j-1728|^2$ leaves that vacuum \emph{massless} (Sec.~\ref{sec:setup}).
(ii) A single extra term, $\propto \ln(1+|j-1728|/1728)$, gives the vacuum a mass \emph{and} tilts the plateau. The tilt raises $n_s$ into the range preferred by CMB+DESI while keeping inflation inside a perturbative string-theory bound, $\mathrm{Im}\,\tau \lesssim 25$ (Secs.~\ref{sec:twoterm}--\ref{sec:predictions}).
(iii) With and without this tilt, the model has two ``faces'': an untilted one (Kairos) and a tilted one (Chronos). We derive both in closed form and find a consistency relation that predicts the running $\alpha_s$ from $(n_s, r)$ with no free parameter (Sec.~\ref{sec:janus}).
(iv) A full consistency chain, covering reheating, supersymmetry breaking, moduli and baryogenesis, passes, and leads to a neutrino-mass prediction, $\sum m_\nu \simeq 0.059$~eV with normal ordering (Secs.~\ref{sec:reheating}--\ref{sec:baryo}).

We are explicit about what is not derived. The size of the tilt parameter, $\varepsilon \approx 10^{-4}$, was chosen after seeing current data, and we have not found its microscopic origin (Sec.~\ref{sec:limits}). Most ingredients are known: the tilt mechanism is that of singular attractors~\cite{arXiv_2512_02969}, and vacua at $\rho$ appear in Ref.~\cite{arXiv_2408_05203}. What is new is the combination, and the chain of checks around it. All predictions are frozen for future data (Sec.~\ref{sec:forward}), and every computation was pre-registered with a hash before it was run.

\section{Modular set-up and what symmetry fixes at $\tau = i$}\label{sec:setup}
We take one complex modulus $\tau = x + i y$ in the upper half plane, with the $\alpha$-attractor K\"ahler potential
\begin{equation}
K = -3\alpha \ln(2\,\mathrm{Im}\,\tau), \qquad \alpha = \tfrac13,
\end{equation}
so the kinetic term is $(\partial x^2 + \partial y^2)/(4y^2)$. Along $x = 0$ the canonical field is $\phi = \ln y/\sqrt{2}$ (reduced Planck units). Following Kallosh and Linde~\cite{arXiv_2408_05203}, the potential is placed in the metric of a nilpotent sector, so any $SL(2,\mathbb{Z})$-invariant function of $\tau$ can serve as $V$.

\paragraph{Stationarity.} The generator $S:\tau\to-1/\tau$ fixes $\tau = i$ and acts there as $\delta\tau \to -\delta\tau$. Hence any $S$-invariant potential that is smooth at $i$ has $\nabla V(i) = 0$. The point $\tau = i$ is therefore always a critical point. This is derived, not chosen. We checked it numerically for every potential used below (D574).

\paragraph{Weight parity.} A modular form of weight $w$ satisfies $f(i) = i^{w} f(i)$, so it vanishes at $i$ unless $w \equiv 0 \pmod 4$. In particular $E_6(i) = 0$ (numerically $|E_6(i)| < 10^{-15}$, with a simple zero), while $E_4(i) = 1.456$. Since $j - 1728 \propto E_6^2/\Delta$, the function $j-1728$ has a \emph{double} zero at $i$ (fitted order $2.000$). A superpotential carrying one factor of $E_6$ therefore has a supersymmetric Minkowski point at $\tau = i$. This is consistent with the no-go result that a de Sitter vacuum cannot sit at $i$ in this minimal setting~\cite{arXiv_2212_03876}. What symmetry does \emph{not} fix is whether $i$ is a minimum or a saddle, nor the value of the vacuum energy.

\paragraph{The plateau.} We build the plateau from the invariant
\begin{equation}
I(\tau) = \frac{1}{4\pi}\ln\!\left(1 + \frac{|j(\tau)-1728|^2}{1728^2}\right), \qquad V_{\rm pl} = V_0\,\frac{I^p}{I^p+1},
\end{equation}
which vanishes only at $i$ (and its images), is single-valued, and grows as $I \simeq y$ toward the cusp. We use $p = 2$. That is the smallest power for which horizon exit occurs at $y_* \lesssim 25$, inside the bound that perturbative heterotic thresholds place on $\mathrm{Im}\,\tau$ (about $k\,\alpha^{-1}(M_s) \approx 29$; D564, D566).

\paragraph{A flaw: the vacuum is massless.} Because $|j - 1728|^2 \propto |\tau - i|^4$, the plateau rises as $V_{\rm pl} \propto |\tau - i|^{4p}$ near the vacuum: quartic for $p=1$ and octic for $p=2$ (fitted exponents $4.00$ and $8.00$). The vacuum has no quadratic term, so no mass. After inflation the field oscillates anharmonically, with equation of state $w = (n-2)/(n+2) = 3/5$ for $n = 8$, and this changes the inflationary e-fold count $N_*$. Earlier predictions that assumed matter-like reheating must be revised (D567a). Section~\ref{sec:twoterm} supplies the missing mass.

%==SEC3==
%==SEC4==
%==SEC5==
%==SEC6==
%==SEC7==
%==SEC8==
%==SEC9==
%==SEC10==
%==SEC11==
%==ACK==
\bibliography{refs}
\end{document}
